\section{Hallazgos de Literatura Adicional}

Esta sección presenta los hallazgos de una revisión exhaustiva de literatura reciente (2020-2025) en temas clave relacionados con la transformación digital en el agro. Se identificaron artículos de fuentes como IEEE, Elsevier y revistas especializadas en agricultura digital, enfocándose en su aplicabilidad al Portal Agro Comercial del Huila.

\subsection{IoT Agrícola}

Farooq et al. (2020) revisan el rol de IoT en agricultura, destacando beneficios como monitoreo remoto y eficiencia en recursos, pero desafíos en conectividad rural. Relacionado con el portal, permite integrar sensores para datos en tiempo real, reduciendo intermediarios al proporcionar información verificable.

Kumar y Sharma (2022) abordan desafíos de conectividad en IoT agrícola, proponiendo soluciones como redes LoRa para áreas remotas como el Huila. Mejora el portal con infraestructura de bajo costo para monitoreo de fincas.

Díaz (2023) analiza plataformas IoT en agricultura, enfatizando integración con big data para predicciones. Aplica al portal como hub central para datos IoT, optimizando decisiones de comercialización.

Gráfica sugerida: Gráfico de líneas mostrando reducción de costos operativos con IoT en agricultura.

\subsection{E-commerce Rural}

Vega (2022) identifica factores de confianza en e-commerce agrícola, como garantías de calidad y pagos seguros. El portal puede implementar certificaciones digitales para construir confianza y eliminar intermediarios.

Herrera (2024) propone soluciones logísticas para e-commerce rural en Colombia, incluyendo alianzas con transportistas locales. Integra en el portal un módulo de logística integrada, facilitando entregas directas.

Gómez et al. (2021) evalúan oportunidades de e-commerce en horticultura colombiana, reduciendo intermediación mediante marketplaces en línea. El portal se basa en este modelo para productos del Huila.

Gráfica sugerida: Diagrama de flujo del proceso de e-commerce rural, desde pedido hasta entrega.

\subsection{Transformación Digital Rural en Colombia y Latam}

Ramírez (2021) discute políticas de inclusión digital en agricultura, recomendando capacitación y acceso a TIC. El portal incluye programas de alfabetización digital para cerrar brechas en el Huila.

Silva y Torres (2023) analizan lecciones de transformación digital rural en Colombia, destacando alianzas público-privadas. Sugiere replicar en el portal con colaboraciones gubernamentales.

Rodríguez (2023) identifica patrones de digitalización en Latinoamérica, enfatizando adopción gradual. Aplica al portal como modelo escalable para la región.

Gráfica sugerida: Mapa regional de niveles de transformación digital en países latinoamericanos.

\subsection{Agricultura de Precisión}

Zhang y Wang (2021) revisan adopción de tecnologías de precisión en países en desarrollo, destacando barreras económicas. El portal ofrece herramientas accesibles, como mapas geoespaciales, para fincas pequeñas del Huila.

Castillo (2023) evalúa avances en agricultura de precisión en Colombia, proponiendo integración con drones y sensores. Mejora el portal con módulos de agricultura 4.0.

Vargas (2024) explora herramientas geoespaciales en agricultura colombiana, permitiendo "cultivar a ciegas" con datos precisos. Integra en el portal para optimizar rendimientos.

Gráfica sugerida: Scatter plot de adopción vs. tamaño de finca en agricultura de precisión.

\subsection{Plataformas Tecnológicas del Sector Agropecuario}

Mendoza (2024) analiza integración de big data en cadenas de suministro agrícolas, mejorando trazabilidad. El portal actúa como plataforma para conectar productores y compradores con datos en tiempo real.

Ojeda-Beltrán (2023) extiende plataformas en agricultura 4.0, enfocando en necesidades reales de productores. Adapta el portal a contextos colombianos.

López (2024) propone plataformas locales para agricultura digital, enfatizando interoperabilidad. Desarrolla el portal como solución regional.

Gráfica sugerida: Diagrama de arquitectura de plataformas tecnológicas agropecuarias.

\subsection{Brecha Digital Rural}

FAO (2021) discute digitalización para agricultura sostenible, identificando brechas en acceso. El portal cierra estas brechas mediante interfaces inclusivas.

CEPAL (2022) analiza brechas digitales en Latinoamérica agrícola, recomendando inversiones en infraestructura. Aplica al Huila con políticas de conectividad.

Deichmann et al. (2016) evalúan si tecnologías digitales transformarán agricultura en desarrollo, concluyendo que sí, con inversiones adecuadas. Soporta la viabilidad del portal.

Gráfica sugerida: Gráfico de barras de brechas digitales por indicadores (acceso, capacitación, etc.).

\subsection{Innovación en Cadenas de Suministro Agrícolas}

Kamilaris et al. (2017) revisan análisis de big data en agricultura, mejorando eficiencia en cadenas. El portal innova con IA para predicciones de demanda.

Fernández (2024) propone integración de sistemas en agricultura 4.0 en Colombia, reduciendo intermediación. Implementa en el portal blockchain para trazabilidad.

Lindgreen y Hingley (2020) discuten marketing sostenible en agricultura, innovando en cadenas de valor. Integra storytelling digital en el portal.

Gráfica sugerida: Diagrama de red de innovación en cadenas de suministro agrícolas.